% GENERATED FILE. Edit canonical lessons, then run npm run generate:books.
% canonical-source-hash: fnv1a64:a3fcfbf86c9184fa
% canonical-lessons: TE-C153-evaru, TE-C153-eppudu, TE-C153-ekkada, TE-C153-evaraina, TE-C153-emaina

\chapter{Nobody, Never, Nowhere, Anyone, Anything}
\label{ch:te-nobody-never-nowhere-anyone-anything}
\hlchaptermodality{153}

\begin{chapteropening}
\textbf{By the end of this chapter:} \emph{I can say what is not so, and ask yes-or-no questions.}

Complete the last lesson of chapter 153: I can say what is not so, and ask yes-or-no questions.
\end{chapteropening}

\section[\texorpdfstring{\te{ఎవరూ} (evarū)}{ఎవరూ (evarū)}]{\te{ఎవరూ} (evarū) — nobody (with a negative verb)}

\label{lesson:TE-C153-evaru}



\begin{quote}
Before the new one: say the Telugu for shade, then the Telugu for a stream.
\end{quote}

\subsection*{You'll want to know: \te{ఎవరూ}}
\textbf{\te{ఎవరూ}} — \emph{evarū} — \textquotedblleft{}nobody (with a negative verb)\textquotedblright{}. Telugu negates a past verb with \textbf{\te{లేదు}} (\emph{lēdu}) joined onto it: \textbf{\te{ఎవరూ} \te{రాలేదు}} — \emph{evarū rālēdu} — \textquotedblleft{}Nobody came.\textquotedblright{} To say something is not a thing, it uses \textbf{\te{కాదు}} (\emph{kādu}): \textbf{\te{ఇది} \te{నా} \te{ఇల్లు} \te{కాదు}} — \emph{idi nā illu kādu} — \textquotedblleft{}This is not my house.\textquotedblright{}

\begin{grammarlens}[title={What you've built}]
One more word for this chapter.
\end{grammarlens}

\subsection*{Your turn}
\begin{itemize}
\raggedright
  \item \emph{Say it:} \emph{evarū}
  \item \emph{Say it:} \emph{evarū}, once more
  \item \emph{Say it:} say \emph{evarū}
  \item \emph{Recall:} say the Telugu for shade, then the Telugu for a stream, then say \emph{evarū} again
\end{itemize}

\begin{tcolorbox}[breakable,colback=teal!4,colframe=teal!35!black,title={Before you move on}]
What does \te{ఎవరూ} mean? (\textquotedblleft{}Nobody (with a negative verb)\textquotedblright{}.) Say it once more.
\end{tcolorbox}
\section[\texorpdfstring{\te{ఎప్పుడూ} (eppuḍū)}{ఎప్పుడూ (eppuḍū)}]{\te{ఎప్పుడూ} (eppuḍū) — never (with a negative verb); always}

\label{lesson:TE-C153-eppudu}



\begin{quote}
Before the new one: say the Telugu for a stream, then the Telugu for nobody.
\end{quote}

\subsection*{You'll want to know: \te{ఎప్పుడూ}}
\textbf{\te{ఎప్పుడూ}} — \emph{eppuḍū} — \textquotedblleft{}never (with a negative verb); always\textquotedblright{}. \textbf{\te{నేను} \te{ఎప్పుడూ} \te{మాంసం} \te{తినను}} — \emph{nēnu eppuḍū māṁsaṁ tinanu} — \textquotedblleft{}I never eat meat.\textquotedblright{}

\begin{grammarlens}[title={What you've built}]
One more word for this chapter.
\end{grammarlens}

\subsection*{Your turn}
\begin{itemize}
\raggedright
  \item \emph{Say it:} \emph{eppuḍū}
  \item \emph{Say it:} \emph{eppuḍū}, once more
  \item \emph{Say it:} say \emph{eppuḍū}
  \item \emph{Recall:} say the Telugu for a stream, then the Telugu for nobody, then say \emph{eppuḍū} again
\end{itemize}

\begin{tcolorbox}[breakable,colback=teal!4,colframe=teal!35!black,title={Before you move on}]
What does \te{ఎప్పుడూ} mean? (\textquotedblleft{}Never (with a negative verb); always\textquotedblright{}.) Say it once more.
\end{tcolorbox}
\section[\texorpdfstring{\te{ఎక్కడా} (ekkaḍā)}{ఎక్కడా (ekkaḍā)}]{\te{ఎక్కడా} (ekkaḍā) — nowhere, not anywhere (with a negative verb)}

\label{lesson:TE-C153-ekkada}



\begin{quote}
Before the new one: say the Telugu for nobody, then the Telugu for never.
\end{quote}

\subsection*{You'll want to know: \te{ఎక్కడా}}
\textbf{\te{ఎక్కడా}} — \emph{ekkaḍā} — \textquotedblleft{}nowhere, not anywhere (with a negative verb)\textquotedblright{}. \textbf{\te{నాకు} \te{ఆ} \te{పుస్తకం} \te{ఎక్కడా} \te{దొరకలేదు}} — \emph{nāku ā pustakaṁ ekkaḍā dorakalēdu} — \textquotedblleft{}I couldn't find that book anywhere.\textquotedblright{}

\begin{grammarlens}[title={What you've built}]
One more word for this chapter.
\end{grammarlens}

\subsection*{Your turn}
\begin{itemize}
\raggedright
  \item \emph{Say it:} \emph{ekkaḍā}
  \item \emph{Say it:} \emph{ekkaḍā}, once more
  \item \emph{Say it:} say \emph{ekkaḍā}
  \item \emph{Recall:} say the Telugu for nobody, then the Telugu for never, then say \emph{ekkaḍā} again
\end{itemize}

\begin{tcolorbox}[breakable,colback=teal!4,colframe=teal!35!black,title={Before you move on}]
What does \te{ఎక్కడా} mean? (\textquotedblleft{}Nowhere, not anywhere (with a negative verb)\textquotedblright{}.) Say it once more.
\end{tcolorbox}
\section[\texorpdfstring{\te{ఎవరైనా} (evarainā)}{ఎవరైనా (evarainā)}]{\te{ఎవరైనా} (evarainā) — anyone, someone}

\label{lesson:TE-C153-evaraina}



\begin{quote}
Before the new one: say the Telugu for never, then the Telugu for nowhere.
\end{quote}

\subsection*{You'll want to know: \te{ఎవరైనా}}
\textbf{\te{ఎవరైనా}} — \emph{evarainā} — \textquotedblleft{}anyone, someone\textquotedblright{}. A yes-or-no question ends in \textbf{-\te{ఆ}} (\emph{-ā}): \textbf{\te{ఉన్నారు}} becomes \textbf{\te{ఉన్నారా}}. \textbf{\te{అక్కడ} \te{ఎవరైనా} \te{ఉన్నారా}?} — \emph{akkaḍa evarainā unnārā?} — \textquotedblleft{}Is anyone there?\textquotedblright{}

\begin{grammarlens}[title={What you've built}]
One more word for this chapter.
\end{grammarlens}

\subsection*{Your turn}
\begin{itemize}
\raggedright
  \item \emph{Say it:} \emph{evarainā}
  \item \emph{Say it:} \emph{evarainā}, once more
  \item \emph{Say it:} say \emph{evarainā}
  \item \emph{Recall:} say the Telugu for never, then the Telugu for nowhere, then say \emph{evarainā} again
\end{itemize}

\begin{tcolorbox}[breakable,colback=teal!4,colframe=teal!35!black,title={Before you move on}]
What does \te{ఎవరైనా} mean? (\textquotedblleft{}Anyone, someone\textquotedblright{}.) Say it once more.
\end{tcolorbox}
\section[\texorpdfstring{\te{ఏమైనా} (ēmainā)}{ఏమైనా (ēmainā)}]{\te{ఏమైనా} (ēmainā) — anything, something}

\label{lesson:TE-C153-emaina}



\begin{quote}
Before the new one: say the Telugu for nowhere, then the Telugu for anyone.
\end{quote}

\subsection*{You'll want to know: \te{ఏమైనా}}
\textbf{\te{ఏమైనా}} — \emph{ēmainā} — \textquotedblleft{}anything, something\textquotedblright{}. \textbf{\te{మీకు} \te{ఏమైనా} \te{కావాలా}?} — \emph{mīku ēmainā kāvālā?} — \textquotedblleft{}Do you need anything?\textquotedblright{} Read two words again: \textbf{\te{ఛాయ}} (\emph{chāya}, shade) is written with \te{ఛ}, and \textbf{\te{ఝరి}} (\emph{jhari}, a stream) with \te{ఝ}.

\begin{grammarlens}[title={What you've built}]
One more word for this chapter.
\end{grammarlens}

\subsection*{Your turn}
\begin{itemize}
\raggedright
  \item \emph{Say it:} \emph{ēmainā}
  \item \emph{Say it:} \emph{ēmainā}, once more
  \item \emph{Say it:} say \emph{ēmainā}
  \item \emph{Recall:} say the Telugu for nowhere, then the Telugu for anyone, then say \emph{ēmainā} again
\end{itemize}

\begin{tcolorbox}[breakable,colback=teal!4,colframe=teal!35!black,title={Before you move on}]
What does \te{ఏమైనా} mean? (\textquotedblleft{}Anything, something\textquotedblright{}.) Say it once more.
\end{tcolorbox}
